\scanpage{31}

\[
\bigl[\,X\ \text{quasi-compact}\Longleftrightarrow
\text{Si } A=\sum_{i\in I}J_i,\ \exists\,J\ \text{fini}\subset I,\ A=\sum_{i\in J}J_i.
\]

\deleted{\unsure{Parties irréductibles}}

Si $J$ est un idéal de $A$, $\mathrm{Spec}\,(A/J)=Y\subset X$ est un sous-espace topologique fermé de $X$. \deleted{On trouve} Il ne change pas si on remplace $J$ par $\widetilde J$, donc \deleted{\unsure{pour étu}} un sous-espace fermé de $X$ est tjrs hom. au spectre d'un anneau réduit.

\uline{Prop} Si $A$ est réduit, $X=\mathrm{Spec}\,A$ irréductible $\Longleftrightarrow A$ intègre.

\[
\begin{aligned}
&\text{En effet, $\mathrm{Spec}\,A$ irréductible}\ \overset{\text{déf}}{\Longleftrightarrow}
\left\{\begin{aligned}
&X\ne\varnothing\\
&\underset{X'\ \mathrm{et}\ X''\ \mathrm{fermés}}{X=X'\cup X''}\Rightarrow X'\text{ ou }X''=X
\end{aligned}\right.\\[-1pt]
&\qquad\begin{array}{l|l}
\text{i.e.}& A\ne0\\
&\text{i.e. }J'\cdot J''=0\Rightarrow J'\text{ ou }J''\ \text{nul}
\end{array}
\end{aligned}
\]

\uline{Corollaire} Soit $A$ quelconque, \deleted{\unsure{il y a correspondance biun|$A$ a une correspondance biunivoque}} les parties fermées irréductibles de $A$ sont les $V(p)$, $p$ idéal premier de $A$ (correspondance biunivoque \unsure{et|$\updownarrow$} renversant l'ordre) et \deleted{\unsure{l'id}} si $x$ correspond à $p$, alors $V(p)$ est égal à $\overline{\{x\}}$. En particulier, $X$ est \uline{sobre}.

\margin{$\mathrm{Spec}(A_f)\simeq X_f$ (homéomorphisme)}

% slanted
\margin{\textcircled{$10'$} Fonctorialité de \unsure{l'esp.} top. $\mathrm{Spec}\,A$ Cas des épi. d'ann.}

\uline{NB} L'espace topologique $\mathrm{Spec}\,A$ ne dépend que de la structure d'anneau de $A$, i.e. de $\mathbb Z$-algèbre \ldots.

\textcircled{11} \uline{Spectre d'une $k$-algèbre : faisceau d'anneaux}

On ne peut réconstituer $A$ à partir du seul espace topologique $X=\mathrm{Spec}\,A$ !

Faisceaux sur $X$ : $\forall f\in A$, ensemble $P(f)$

$\forall f,g\in A$ tels que $\deleted{\unsure{D}}X_{\unsure{fS|f}}\subset \deleted{\unsure{D}}X_g)$, une application $P(g)\to P(f)$

\[
\begin{aligned}
&\text{i.e. }V(f)\supset V(g)\\
&\text{i.e. }\widetilde{fA}\subset\widetilde{gA}\\
&\text{i.e. }f\in\widetilde{gA}\\
&\text{i.e. }\left\{\begin{aligned}&\exists\,n\in\mathbb N,\ f^n=gh\\&h\in A\end{aligned}\right.\ \big|\ \text{écrivons } f\prec g
\end{aligned}
\]

avec transitivité si $f,g,h$ tels que $f\prec g\prec h$ et identité si $f=g$

Avec $\forall f\in A$ et $(f_i)_{i\in I}$, $f_i\in{I}$ tels que

\[
\begin{aligned}
&X_f=\bigcup X_{f_i}\\
&\text{i.e. }V(f)=\bigcap V(f_i)\\
&\text{i.e. }\widetilde{fA}=\widetilde{\sum f_i A}\\
&\text{i.e. }\left\{\begin{aligned}
&\forall i\quad f_i\prec f\quad\unsure{\vee|\checkmark}\\
&\exists n\in\mathbb N\text{ et }(h_i)_{i\in I}\text{ avec } f^n=\sum f_i h_i
\end{aligned}\right.
\end{aligned}
\]

on veut exactitude dans

\[
P(f)\longrightarrow\prod_i P(f_i)\rightrightarrows\prod_{ij}P(f_i f_j)
\]

% archive sheet number: 30
